\documentclass[11pt]{article}
\usepackage[margin=1in]{geometry}
\usepackage{amsmath,amssymb}
\usepackage{booktabs}
\usepackage{array}
\usepackage[hidelinks]{hyperref}
\usepackage{enumitem}

\title{\textbf{Comparing Exploration Algorithms on Ms.\ Pac-Man}\\
\large A Course-Project Experiment Design\\
\normalsize \emph{Challenging Problems in Reinforcement Learning} (SS~2026)\\
\normalsize Topic: Exploration vs.\ Exploitation}
\author{}
\date{2026-05-31}

\newcommand{\score}{S}

\begin{document}
\maketitle

\begin{abstract}
\noindent We compare six exploration mechanisms---classical action-selection rules and deep
intrinsic-motivation methods, including Learning Progress Monitoring (LPM)---on the Atari
\textsc{Ms.\ Pac-Man} environment, measure their relative strength, and analyse \emph{why} they
differ. The study is built on the LPM codebase of Hou, An \& Du (2026), \emph{Beyond Noisy-TVs:
Noise-Robust Exploration via Learning Progress Monitoring}. Two questions drive the design:
\textbf{(RQ1)} how robust is each method to an injected noisy-TV, and \textbf{(RQ2)} does LPM's
advantage on \emph{clean} Ms.\ Pac-Man grow with environment stochasticity---a probe of the
hypothesis that the game's random ghost movement already acts as aleatoric noise.
All methods share one PPO base so that observed differences are attributable to the exploration
mechanism rather than the base learner.
\end{abstract}

\section{Motivation and research questions}
The course topic is the exploration--exploitation trade-off. Classical exploration
(\(\varepsilon\)-greedy, softmax/Boltzmann, UCB, Thompson sampling) was designed for bandit and
value-based settings; deep RL instead uses intrinsic-motivation bonuses (prediction error, novelty,
learning progress). The \emph{noisy-TV} problem is the canonical failure mode of prediction-error
exploration: an unlearnable source of randomness produces perpetually high prediction error, so a
curiosity-driven agent fixates on it. LPM rewards the \emph{improvement} of its dynamics model
rather than the error level, and is therefore noise-robust by construction.

\begin{description}[leftmargin=2.2em,style=nextline]
  \item[RQ1 (noise robustness, reproduction).] Under an action-triggered noisy-TV, which methods
  collapse and which survive? How large is LPM's advantage over prediction-error curiosity (RND,
  ICM) and over a method explicitly built to filter aleatoric noise (AMA)?
  \item[RQ2 (stochasticity attribution, novelty).] LPM appears to win even on \emph{clean}
  Ms.\ Pac-Man. Hypothesis: the game's random ghost movement acts as aleatoric noise that fools
  prediction-error methods. Because the ghost AI is fixed in the ROM and cannot be tuned, we probe
  the \emph{general} claim---``LPM's advantage grows with environment stochasticity''---using ALE
  sticky actions (\texttt{repeat\_action\_probability}) as a controllable stochasticity knob.
\end{description}

\section{Methods}
All six trained methods use the same PPO base. Intrinsic-motivation methods add an intrinsic reward
with weight \(\beta\); the per-step training reward is
\begin{equation}
  r_t \;=\; r_t^{\,e} \;+\; \beta\, r_t^{\,i},
  \label{eq:reward}
\end{equation}
with \(\beta=1\) throughout. The two classical baselines run on the same PPO base: plain PPO is the
entropy-regularised stochastic policy (i.e.\ softmax/Boltzmann exploration), and
\(\varepsilon\)-greedy adds a uniform-random action layer with probability \(\varepsilon=0.1\).

\begin{table}[h]
\centering
\begin{tabular}{@{}llp{6.4cm}@{}}
\toprule
\textbf{Class} & \textbf{Method} & \textbf{Mechanism (intrinsic reward \(r^{i}\) or rule)} \\
\midrule
Learning progress        & \textbf{LPM} & improvement of the dynamics model, \(\mathbb{E}\!\left[\varepsilon^{(\tau-1)}-\varepsilon^{(\tau)}\right]\) \\
Prediction error         & RND          & error vs.\ a fixed random target network \\
Prediction error (feat.) & ICM          & inverse-dynamics feature prediction error \\
Noise-robust curiosity   & AMA          & curiosity with aleatoric-uncertainty filtering \\
Classical (softmax)      & PPO          & entropy-regularised stochastic policy (no \(r^{i}\)) \\
Classical (\(\varepsilon\)-greedy) & PPO + \(\varepsilon\) & random action w.p.\ \(\varepsilon=0.1\) \\
\bottomrule
\end{tabular}
\caption{The six trained methods. UCB and Thompson sampling are \emph{discussion-only}: neither has
a scalable deep-RL drop-in (UCB needs visit counts/uncertainty; Thompson needs posterior sampling),
which itself illustrates why count-based bonuses and LPM are their practical descendants.}
\end{table}

\section{Experimental design}
Two manipulation axes are kept separate to bound the run count.

\paragraph{RQ1 --- noise axis.} Each method is run in the \emph{clean} game and the \emph{noisy}
variant (action-noise: two idle actions replace the observation with a random CIFAR-10 image),
\(\times\,3\) seeds: \(6\times2\times3 = 36\) runs.

\paragraph{RQ2 --- stochasticity sweep (clean only).} ALE sticky-action probability
\(p \in \{0.0, 0.25, 0.5\}\); the \(p=0.0\) point reuses the RQ1 clean runs. The sweep covers the
four most diagnostic methods (LPM, RND, ICM, plain PPO) \(\times\,\{0.25,0.5\}\,\times\,3\) seeds
\(=24\) runs.

\paragraph{Budget.} \(\approx 60\) runs \(\times \sim\!20\) min on Apple-Silicon MPS
\(\approx 20\) h. Each run trains \(10^{6}\) aggregate frames, 16 parallel environments, PPO with
learning rate \(10^{-4}\), clip \(0.1\), 3 epochs, 8 minibatches, \(\gamma=0.99\), GAE.

\section{Metrics}
Let \(\score\) denote the final extrinsic game score (the mean episode return over the last
\(10\%\) of updates), reported as mean\(\,\pm\,\)std over seeds.

\begin{itemize}[leftmargin=1.4em]
  \item \textbf{RQ1 headline --- noise-robustness drop:}
  \begin{equation}
    \mathrm{Drop} \;=\; \frac{\score_{\mathrm{clean}} - \score_{\mathrm{noisy}}}{\score_{\mathrm{clean}}}.
  \end{equation}
  Prediction: LPM (and AMA) smallest; RND/ICM largest. Reproduces the paper's central result.
  \item \textbf{RQ2 headline --- stochasticity-advantage curve:}
  \(\;\score_{\mathrm{LPM}}(p) - \score_{\mathrm{RND}}(p)\) and the same vs.\ ICM, as functions of
  the sticky probability \(p\). Prediction: increasing \(\Rightarrow\) supports the hypothesis.
  \item \textbf{Attribution (auxiliary):} learning curves (\(\score\) vs.\ frames), and the
  per-method intrinsic-reward trajectory. LPM's \(r^{i}\) should decay toward \(0\) as the world
  model converges; RND/ICM under noise should stay elevated---evidence of being ``stuck on noise.''
  Intrinsic-reward magnitudes are \emph{not} comparable across methods (different scales), so only
  within-method trends are interpreted.
\end{itemize}

\section{Infrastructure and reproducibility}
The embedded LPM Atari snapshot did not run out of the box; eight import/shape/correctness fixes
were required (channel-first frames, ALE registration in subprocess workers, enabling the intrinsic
reward, removing a Sigmoid decoder incompatible with normalised targets, MPS device support; all
logged in \texttt{LPM\_exploration/UPSTREAM.md}). Each run writes a per-update CSV
(\texttt{update, frames, fps, ep\_score\_mean, ep\_score\_std, ep\_len\_mean, int\_rew, ext\_rew,
pred\_loss, unc\_loss, dist\_entropy, value\_loss, action\_loss}). A grid runner launches all runs
sequentially on MPS (resumable: completed CSVs are skipped); an analysis script aggregates the CSVs
into the results table and three figures.

\section{Limitations (stated up front)}
\begin{enumerate}[leftmargin=1.6em]
  \item \textbf{Sticky actions \(\neq\) ghost randomness.} Sticky actions are a controllable
  \emph{proxy} for aleatoric stochasticity, not an isolation of ghost behaviour. RQ2 therefore tests
  the general ``stochasticity \(\to\) LPM edge'' claim, with ghost randomness as qualitative
  motivation only.
  \item \textbf{Short budget, small sample.} \(10^{6}\) frames \(\times\,3\) seeds yields
  \emph{qualitative trends}, not rigorous statistical significance. (At this scale a human casually
  reaches \(\sim\!7000\) while these agents reach \(\sim\!500\); maximising score is not the goal.)
  \item \textbf{No per-state localisation.} We log aggregate intrinsic reward, not its distribution
  over game states, so the noisy-TV mechanism is shown indirectly.
  \item \textbf{Episodic-bonus family absent.} EME and TDD/EDT model files are missing from the
  snapshot, so the comparison spans the curiosity, learning-progress, and classical families, not
  episodic bonuses.
\end{enumerate}

\end{document}
